\section{Why Society is so Toxic}

\paragraph{Old Books}

I like to maintain a library of old books. There is great pleasure in rethinking forgotten thoughts and ideas. Presentism is so tempting because what happens now is apparently more interesting than what happened before. But not to know the past is like believing a fully formed adult never went through childhood. If old ideas are explicitly forgotten, they nevertheless act on us implicitly and unconsciously.

\paragraph{The Cosmic Race}


One such book is \emph{East and West: An Inquiry Concerning World Understanding}, by \textbf{F. S. C. Northrop}, published in 1947. He takes us on a world historical tour of the best in the cultures of Mexico, the United States, British Democracy, German Idealism, Russian Communism, Medieval Catholicism, and finally several cultures of the Orient.

To all appearances, Mexico was deeply Catholic, so Northrop was surprised at the actual diversity of ideas there. Mexico was not immune to European intellectual movements and for a while, was under the sway of the anti-religious ideology of Positivism. Northrop was surprised to find, at the Fiesta Courtyard of the Secretariat of Public Education in Mexico City, sculptures of four figures representing different civilizational influences.

\begin{itemize}
\item \textbf{Quetzalcoatl}: Indigenous Mesoamerica culture
\item \textbf{Bartolomé de las Casas}: Spanish culture
\item \textbf{Plato}: Greek culture
\item \textbf{Buddha}: Indian culture
\end{itemize}


The design was inspired by the writings of the philosopher José Vasconcelos, who was influential at the time. Vasconcelos's goal was the creation of a native Hispanic-American culture. Along with that, in his book \emph{The Cosmic Race}, which promoted the idea of a ``fifth race'', an agglomeration of all the races in the world, as the model for a new civilization. That would be the counterpart to his four civilizational influences.

Ultimately, however, he embraced a deeply Catholic political conservatism. Among others, we can include Ernst Junger, E. F. Schumacher, and Charles Maurras who also came to that vision late in life.

\paragraph{Medieval Catholicism}


Northrup then wrote a thoughtful chapter on the foundations of Medieval Catholicism and its Aristotelean influence. This is not the place for a full description of the chapter, but Northrup shows how the notion of a final formal cause serves as the foundation for the Medieval outlook. It answers:

\begin{itemize}
\item The nature of the human soul
\item The nature of God
\item The nature of divine revelation
\end{itemize}


He points out that the meaning of words like ``soul'', ``God'', ``grace'', or ``revelation'' are not immediately obvious. They need to be worked out logically. It is the genius of the Scholastics to do it so comprehensively.

Northrup then points out how Galileo eventually undercut that grand synthesis. His discovery of the laws of motion showed the inadequacy of Aristotle's understanding of motion. It came at a high cost because Galileo eliminated secondary qualities from physics. Whence, then, are the colors, sounds, tastes, etc.? Although Galileo may be correct regarding the physical and corporeal states, there is no place in his system for the subtle states of the human being.

So what if Galileo and describe the movements of the planets or a hit baseball. Can he account for the movements in a small city? Or a family gathering? Or even your own movements throughout the day? There is no formula to predict where you will be three hours from now or what your posture will be at the time.

Can Galileo, for all his admitted genius, tell you the purpose of life or who you should marry? Of course not. Over the centuries we have abandoned an architectonic view of the world to a constricted view, all while considering it progress. No wonder that so many thinking men have come to that conclusion after first having spent many decades thinking ordinary thoughts.

\paragraph{Oil of Vitriol}


As a teenager, I worked during the summers at my father's chemical laboratory. I developed an odd fascination with sulfuric acid. On the one hand, it was dangerous to the touch, but, on the other, it was an excellent cleaning agent. Glass beakers encrusted with chemical residues from resin manufacturing, suddenly became clean in a bath of sulfuric acid.

In our time, old books are destroyed by the vitriol of critical theory. The ideas, knowledge, and insight in such books no longer have value on their own. They are merely masks to hide some sort of oppression.

\paragraph{Subconscious Guilt}


Although there is actually nothing explicitly religious in this lecture by Bishop Fulton Sheen, it is simply assumed. This everyday world of experience is founded on guilt. Not many people want to hear that, but only psychopaths do not experience guilt.

It is best to be consciously aware of your sense of guilt. Then you can take care of and make amends. The act that created the guilt can even be forgiven, restoring some small sense of justice to the cosmos. But unconscious guilt is insidious. Eventually it will erupt in surprising and toxic ways.

Bishop Sheen shows that unconscious guilt manifests in three ways:

\begin{itemize}
\item \textbf{Neurosis}: This includes symptoms like anxiety, sadness or depression, anger, irritability, mental confusion, low sense of self-worth, etc. In the USA, the number of people on psychotropic drugs is alarming. Their society is not producing psychologically healthy people.

\item \textbf{Dreams}: These are more personal. However, the symbols found in dreams can also appear in social contexts. We are awash in symbols that do not mean much.

\item \textbf{Scapegoating}: This is blaming others for one's psychological deficiencies. If the scapegoater feels bad, it must be because someone else is causing it.
\end{itemize}

\url{https://www.youtube.com/watch?v=v32FELEjEG4}

\flrightit{Posted on 2022-09-11 by Cologero}

\begin{center}* * *\end{center}

\begin{footnotesize}
\begin{sffamily}
\texttt{James on 2022-10-14 at 03:55 said:}

I do believe the successors of Galileo or at least those who pay homage to his altar will indeed find a way to express a formula of human motion or human marriage . In some ways , public policy already takes advantage of human flow in cities and meddling with birth rates . However , it is like the advent of so called Artificial Intelligence . It is not that machines are becoming intelligent , it is that humans are becoming more mechanical thus erasing the barrier between human and machine intelligence . Analogously , the formulae for man's marriages and movements will be achieved not because material science has understood humanity , but that humanity has become more material than man .


\hfill

\end{sffamily}
\end{footnotesize}

